\section{机器人基座模型是不是独立学科：So far, not yet}

上一章讲两次范式转移，本章讨论一个争议：机器人基座大模型是否应被视作独立于大模型的学科？谭杰的回答很克制：so far, not yet。当前多数机器人智能仍依赖多模态大模型，只是在其上补足 robot action 输出、动作数据和阶段性 fine-tuning。未来如果遇到新的 data format、世界模型或控制瓶颈，它可能变成更独立的学科，但现在还没有质变。

\begin{figure}[H]
\centering
\includegraphics[width=0.96\textwidth]{vla-action-gap.png}
\caption{VLA 补上 Action 输出：多模态模型要控制机器人，必须学会输出动作。自制概念图，依据 00:13:06--00:23:44 对谈内容整理。}
\end{figure}

\begin{knowledgebox}{读图：从 VLM 到 VLA，关键是动作数据}
Vision 和 Language 进入多模态模型，原本输出文本和计划；机器人需要的是动作，因此要加入 Action Data，让模型从 VLM 变成 VLA。这里的动作不是抽象文字，而是机器人关节、轨迹、末端执行器和控制信号。
\end{knowledgebox}

\subsection{VLA、泛化和成功率}

本节把“动作输出”落到可用性。VLA 的意义不只是把 action 加进名字，而是让机器人从语言和视觉真正走向可执行动作。

VLA 是 Vision-Language-Action，即视觉、语言和动作模型。它希望把“看见环境”“理解指令”和“输出动作”放进同一个模型或系统中。当前模型在简单 pick-and-place 等任务上成功率很高，但对精细操作，例如拉拉链、精确抓取和方向控制，成功率可能只有三四成。这个数字在研究视频里有进展意义，但在现实生活里不可用。

\begin{warningbox}{研究成功率与产品可用性不同}
在真实世界里，30\%--40\% 成功率的精细操作几乎不可用。用户需要的是稳定、可恢复、可解释和可安全失败的系统，而不只是一次视频里完成任务。
\end{warningbox}

\subsection{从 idea 到 demo 到落地}

谭杰用自动驾驶做类比：从一个想法到论文 prototype 可能只需半年到一年；从论文到敢做 live demo 可能要一两年；从 live demo 到真正落地可能要五到十年。机器人比自动驾驶更难，因为动作空间更大、任务更多、物理交互更复杂。自动驾驶是动作输出相对有限的垂直场景，仍然花了十多年才接近落地。

\begin{knowledgebox}{从研究到落地的阶段表}
\begin{tabularx}{\textwidth}{p{0.22\textwidth}p{0.32\textwidth}X}
\toprule
阶段 & 常见产物 & 主要风险 \\
\midrule
Idea & 研究想法、训练 recipe、初步假设 & 可能只在很窄条件下成立。 \\
Prototype & 论文、离线评测、录制视频 & 假设很多，失败样本被过滤。 \\
Live demo & 现场演示、有限任务成功 & 稳定性和异常恢复仍不足。 \\
Pilot & 小规模客户或场景试点 & 成本、维护、安全和运营压力出现。 \\
Production & 可复制商业部署 & 需要可靠硬件、数据飞轮和服务体系。 \\
\bottomrule
\end{tabularx}
\end{knowledgebox}

\begin{importantbox}{机器人时间尺度}
研究突破和商业落地不是同一件事。机器人需要从任务成功率、硬件可靠性、安全、成本、维护、数据飞轮和场景价值上同时过线，才能从 demo 进入真实生产。
\end{importantbox}

\subsection{本章小结}

机器人基座模型目前更像多模态大模型的延伸，而不是完全独立学科。真正决定它是否独立的，是未来是否需要新的数据格式、世界模型、控制范式和跨本体泛化机制。

